% Plan del Taller virtual ATA EPE: simulador físico para armar con las piezas del kit.
% Compilar:  latexmk -pdf plan_taller.tex
\documentclass[11pt,a4paper]{article}
\usepackage[spanish,es-noshorthands,es-tabla]{babel}
\usepackage[utf8]{inputenc}
\usepackage[T1]{fontenc}
\usepackage{lmodern}
\usepackage[margin=2.2cm]{geometry}
\usepackage[table]{xcolor}
\usepackage{amsmath,amssymb,bm}
\usepackage{tabularx,array,booktabs}
\usepackage{enumitem}
\usepackage{fontawesome5}
\usepackage[most]{tcolorbox}
\usepackage[hidelinks]{hyperref}

\definecolor{morado}{HTML}{2A1B7E}
\definecolor{amarillo}{HTML}{FFCF2E}
\definecolor{rojo}{HTML}{C4142C}
\definecolor{verde}{HTML}{1F9D55}
\definecolor{gris}{HTML}{5D5A73}

\newcolumntype{Y}{>{\raggedright\arraybackslash}X}
\renewcommand{\arraystretch}{1.3}
\setlist{itemsep=1pt,topsep=3pt}
\numberwithin{equation}{section}

\newcommand{\vect}[1]{\bm{#1}}
\newcommand{\R}{\mathbf{R}}
\newcommand{\I}{\mathbf{I}}
\newcommand{\E}{\mathbb{1}}
\newcommand{\skw}[1]{[\vect{#1}]_\times}

\newtcolorbox{decision}[1]{colback=amarillo!12,colframe=amarillo!80!black,fonttitle=\bfseries,title={\faLightbulb~#1}}
\newtcolorbox{prueba}[1]{colback=verde!7,colframe=verde,fonttitle=\bfseries,title={\faCheck~#1}}

\title{\textcolor{morado}{\textbf{Taller virtual ATA EPE}}\\[4pt]
  \large Un simulador físico del kit construido desde cero:\\ proyecto de estudio en álgebra geométrica,
  mecánica analítica y programación}
\author{GIVIE · ATA EPE Robótica Escolar}
\date{29 de octubre de 2026}

\begin{document}
\maketitle

\begin{tcolorbox}[colback=morado,colframe=morado,coltext=white,arc=3mm]
\textbf{\textcolor{amarillo}{La idea en una frase.}} Una mesa de trabajo en el navegador donde se toman vigas,
ruedas, tornillos, tuercas y bandas elásticas, se unen como en la vida real, y el conjunto obedece a la física:
gravedad, fricción, inercia, elasticidad. Si se aprieta la tuerca, las piezas se vuelven un solo cuerpo; si se deja
floja, giran; si se gira la manivela, todo lo que está conectado responde.
\end{tcolorbox}

\begin{tcolorbox}[colback=amarillo!15,colframe=amarillo!80!black,arc=3mm]
\textbf{Dos caras de un mismo proyecto.} \emph{Hacia el público}: un taller virtual donde cualquiera arma con el kit.
\emph{Hacia adentro}: un proyecto de estudio donde cada pieza del simulador se construye desde cero, después de
entender la física que la sustenta, su formulación en álgebra geométrica proyectiva (PGA) y en mecánica analítica, y
su discretización. Nada se programa sin haberse derivado antes.
\end{tcolorbox}

\tableofcontents
\newpage

% ==========================================================================================
\section{Objetivos y alcance}

\subsection{Lo que se podrá hacer}
\begin{enumerate}[leftmargin=*]
  \item \textbf{Armar}: tomar piezas del catálogo (las mismas de la web), moverlas, girarlas y encajarlas en la
        retícula de 10\,mm del kit.
  \item \textbf{Unir}: pasar un \emph{tornillo}, que es una pieza aparte, por agujeros alineados de dos o más piezas,
        y ponerle o no una \emph{tuerca}, fija o floja.
  \item \textbf{Simular}: con gravedad siempre presente, ver caer, balancear, girar y transmitir movimiento.
  \item \textbf{Transmitir}: ruedas que se tocan, de lado o de forma transversal, y bandas elásticas entre poleas.
  \item \textbf{Actuar}: girar una manivela, empujar, arrastrar o soltar piezas.
  \item \textbf{Medir y explicar}: velocidades angulares, relaciones de transmisión, tensiones, energía; botón
        «Ver las cuentas», como en la balanza de la web.
  \item \textbf{Guardar y compartir} un montaje, y llevarlo a realidad aumentada.
\end{enumerate}

\subsection{Lo que no se hará (por ahora)}
Deformación de vigas y pletinas (se tratan como sólidos rígidos), electrónica y motores eléctricos (fase
posterior, para los robots electromecánicos), y fluidos.

\begin{decision}{Principio de diseño}
La física no es decorativa. Cada fenómeno tiene un modelo explícito (este documento), parámetros con unidades que
se pueden \emph{medir en el aula} con el kit real, y una prueba con resultado analítico conocido. Así el simulador
sirve para enseñar: lo que se ve en pantalla se puede explicar con las mismas ecuaciones.
\end{decision}

% ==========================================================================================
\section{Método de trabajo: estudiar, discretizar, programar}\label{sec:metodo}

Cada módulo del simulador (sección~\ref{sec:modulos}) sigue el mismo ciclo de cinco pasos:

\begin{enumerate}[leftmargin=*]
  \item \textbf{Estudiar.} Una \emph{ficha de estudio} en LaTeX con la teoría del módulo, las lecturas precisas
        (capítulos y secciones de la bibliografía) y ejercicios a mano. El módulo no avanza hasta poder responder sus
        \emph{preguntas de control}.
  \item \textbf{Discretizar.} Pasar de las ecuaciones continuas a un algoritmo: qué se guarda en memoria, qué se
        aproxima, con qué error y con qué estabilidad. Se escribe en la misma ficha, en pseudocódigo.
  \item \textbf{Probar primero.} Antes de programar, se escriben las pruebas: identidades algebraicas y resultados
        analíticos que el código debe reproducir (sección~\ref{sec:verificacion}).
  \item \textbf{Programar.} Se implementa en JavaScript, sin librerías de física. Cada función lleva en su comentario
        la ecuación que implementa.
  \item \textbf{Publicar.} Lo construido se convierte en una función visible para el público, con su «Ver las cuentas».
\end{enumerate}

\begin{decision}{Qué se construye desde cero y qué no}
\textbf{Desde cero}: el álgebra (PGA), la cinemática, la dinámica de cuerpo rígido, los integradores, las ligaduras,
la detección de colisiones, el contacto con fricción, las transmisiones del kit y el análisis del mecanismo.
\textbf{Se reutiliza}: solo el dibujo (three.js). Pero sus entradas y salidas se reinterpretan: el motor físico
entrega \emph{motores} de PGA, que un adaptador convierte en matrices para dibujar; el cursor y el dedo llegan como
\emph{rectas} de PGA (el rayo desde la cámara), con las que se eligen y mueven piezas.
\end{decision}

% ==========================================================================================
\section{Arquitectura}

\subsection{Tecnologías}
\begin{tabularx}{\linewidth}{@{}l Y@{}}
\toprule
\textbf{Capa} & \textbf{Elección y motivo}\\
\midrule
Plataforma & Página web estática, igual que el sitio: funciona en computador, tableta y celular, sin instalar nada.\\
Gráficos & three.js 0.169 (ya usado), con los modelos GLB de las piezas que ya existen.\\
Núcleo físico & Propio, en JavaScript y formulado en PGA: cuerpos rígidos, integradores, ligaduras, colisiones y
  contacto con fricción.\\
Capa del kit & Grafo de ensamble, tornillos y tuercas, cuerpos compuestos, ruedas de fricción y bandas
  elásticas.\\
Análisis & Grados de libertad, ejes, coordenadas generalizadas, lagrangiano reducido y trabajo virtual.\\
Pruebas & \texttt{node --test} (incluido en Node 22) y navegador sin interfaz para las pruebas visuales.\\
Interfaz & HTML y CSS con el estilo de la web; controles táctiles.\\
\bottomrule
\end{tabularx}

\begin{decision}{¿Por qué un motor propio?}
Porque el objetivo no es solo tener un simulador, sino entender cada línea. Un motor existente (como Rapier) sería
más rápido de integrar, pero sería una caja negra. El precio es tiempo: el contacto con fricción es la parte más
difícil y se le reserva un módulo completo. Como respaldo, si un módulo se atasca, se puede usar temporalmente la
detección de colisiones de una librería y seguir con lo demás.
\end{decision}

\subsection{Carpetas del código}
\begin{tabularx}{\linewidth}{@{}l Y@{}}
\toprule
\texttt{nucleo/pga.js} & Álgebra geométrica proyectiva: multivectores, productos, motores, exponencial y logaritmo.\\
\texttt{nucleo/cuerpo.js} & Cuerpo rígido: motor, velocidad, inercia, fuerzas (forques).\\
\texttt{nucleo/integrador.js} & Avance en el tiempo sobre el grupo de movimientos rígidos.\\
\texttt{nucleo/ligaduras.js} & Uniones (revolución, cilíndrica, fija) y su solución.\\
\texttt{nucleo/colision.js} & Formas de colisión y detección (SAT, GJK).\\
\texttt{nucleo/contacto.js} & Contacto unilateral con fricción de Coulomb.\\
\texttt{kit/} & Catálogo con puertos, ensamble, tornillos y tuercas, cuerpos compuestos, ruedas y bandas.\\
\texttt{analisis/} & Grados de libertad, ejes, lagrangiano reducido, trabajo virtual.\\
\texttt{app/} & La cara pública: escena, interfaz, instrumentos, adaptador PGA $\leftrightarrow$ three.js.\\
\texttt{pruebas/} & Pruebas automáticas de cada módulo.\\
\texttt{estudio/} & Fichas de estudio (LaTeX), una por módulo.\\
\bottomrule
\end{tabularx}

\subsection{Ciclo de cada cuadro}
\begin{equation*}
\text{entrada} \;\to\; \text{ensamble (si cambió)} \;\to\; \text{cuerpos} \;\to\;
\underbrace{\text{fuerzas} \to \text{velocidades} \to \text{ligaduras y contactos} \to \text{motores}}_{n \text{ subpasos}} \;\to\; \text{instrumentos} \;\to\; \text{dibujo}
\end{equation*}

% ==========================================================================================
\section{Notación, unidades y piezas}

\subsection{Notación}
Vectores en negrita ($\vect{x}$), matrices en mayúscula recta ($\R$, $\I$). Cada cuerpo $B$ tiene posición del centro
de masa $\vect{x}_B$, orientación $\R_B\in SO(3)$ (en el código, un cuaternión $q_B$), velocidad $\vect{v}_B$ y
velocidad angular $\vect{\omega}_B$. Un punto con coordenadas locales $\vect{r}$ está en
$\vect{p} = \vect{x}_B + \R_B\vect{r}$ y se mueve con $\dot{\vect{p}} = \vect{v}_B + \vect{\omega}_B\times\R_B\vect{r}$.
$\skw{a}$ es la matriz tal que $\skw{a}\vect{b} = \vect{a}\times\vect{b}$.

\subsection{Unidades y escala}
Las piezas miden milímetros y pesan gramos. Los motores físicos tienen tolerancias pensadas para metros, así que se
simula en \textbf{centímetros}, gramos y segundos (CGS):
\begin{equation}
  g = 981\ \mathrm{cm\,s^{-2}},\qquad 1\ \mathrm{N} = 10^{5}\ \mathrm{g\,cm\,s^{-2}},\qquad
  1\ \mathrm{N\,m} = 10^{7}\ \mathrm{g\,cm^{2}\,s^{-2}},
\end{equation}
para que las magnitudes típicas sean de orden 1 y los cálculos estén bien condicionados. Toda magnitud que se
muestra en pantalla se convierte a SI.
La gravedad es $\vect{g} = (0,-g,0)$, \textbf{siempre activa}.

\subsection{Pieza}\label{sec:piezas}
Una pieza $P$ es un sólido rígido con:
\begin{itemize}[leftmargin=*]
  \item Malla visual (los GLB actuales) y forma de colisión: unión de primitivas (cajas, cilindros) o envolventes
        convexas; nunca la malla con agujeros, que es costosa e inestable.
  \item Material $\mathcal{M}$ (plástico, metal, caucho) con densidad $\rho$, coeficientes de fricción y restitución.
  \item Masa $m = \rho V$, centro de masa $\vect{c}$ e inercia $\I$, calculados de la malla cerrada (la herramienta
        \texttt{manifold3d} da $V$ exacto). En un cilindro sólido, por ejemplo,
        $I_{\text{eje}} = \tfrac12 m R^2$ e $I_{\perp} = \tfrac1{12} m(3R^2 + h^2)$.
  \item Un conjunto de \textbf{puertos} $\{(\vect{o}_k, \vect{a}_k, d_k, \ell_k)\}$: cada agujero tiene centro
        $\vect{o}_k$, eje unitario $\vect{a}_k$, diámetro $d_k = 4{,}6$\,mm y largo $\ell_k$. Los tornillos tienen un
        puerto macho de diámetro $4{,}0$\,mm, y las tuercas uno roscado.
\end{itemize}

\begin{decision}{Medir en lugar de suponer}
La densidad real del plástico, la masa de cada pieza y los coeficientes de fricción se miden con el kit, con una
báscula de cocina y un plano inclinado (sección~\ref{sec:calibracion}). Mientras tanto se usan valores típicos: plástico
$\rho\approx1{,}05\ \mathrm{g/cm^3}$, acero $\rho = 7{,}85\ \mathrm{g/cm^3}$.
\end{decision}

% ==========================================================================================
\section{Geometría y cinemática en álgebra geométrica proyectiva}\label{sec:pga}

\subsection{El álgebra}
Se usa el álgebra de Clifford $\mathbb{R}_{3,0,1}$, generada por $e_0, e_1, e_2, e_3$ con
\begin{equation}
  e_0^2 = 0,\qquad e_1^2 = e_2^2 = e_3^2 = 1,\qquad e_ie_j = -e_je_i\ (i\neq j).
\end{equation}
Un multivector tiene $2^4 = 16$ componentes, repartidas en grados 0 a 4. El producto geométrico $ab$ se descompone en
partes de grado definido; de él salen el producto exterior $\wedge$ (unir subespacios), el interior $\cdot$ (proyectar,
medir ángulos), el regresivo $\vee$ (\emph{join} de elementos) y el conmutador $a\times b = \tfrac12(ab - ba)$. La
pseudoescalar es $\mathbf{I} = e_{0123}$.

\subsection{Elementos geométricos}
\begin{tabularx}{\linewidth}{@{}l l Y@{}}
\toprule
\textbf{Objeto} & \textbf{Grado} & \textbf{Forma}\\
\midrule
Plano $ax + by + cz + d = 0$ & 1 & $a e_1 + b e_2 + c e_3 + d e_0$\\
Recta (eje, agujero) & 2 & bivector: parte euclídea (dirección) más parte ideal (momento); coordenadas de Plücker\\
Punto $(x,y,z)$ & 3 & $e_{123} + x\,e_{032} + y\,e_{013} + z\,e_{021}$\\
\bottomrule
\end{tabularx}

\medskip
La recta que pasa por dos puntos es $P\vee Q$; la intersección de dos planos es $p\wedge q$. Los \emph{puertos} de las
piezas (sección~\ref{sec:piezas}) se representan como rectas: un agujero \emph{es} su eje.

\subsection{Movimientos rígidos: motores}
Todo desplazamiento rígido es un \emph{motor} $M$ (un elemento del subálgebra par con $M\widetilde{M} = 1$), que actúa
sobre cualquier elemento por el producto «sándwich»
\begin{equation}
  X' = M\,X\,\widetilde{M},
\end{equation}
y dos movimientos seguidos se componen multiplicando: $M_{21} = M_2M_1$. Los cuaterniones son exactamente los motores
sin traslación. Por el teorema de Chasles, todo motor es un movimiento de \textbf{tornillo}: un giro $\theta$
alrededor de una recta $L$ (normalizada, $L^2 = -1$) combinado con un avance $d$ a lo largo de ella,
\begin{equation}
  M = \exp\!\Bigl(-\tfrac12\,(\theta + d\,\mathbf{I})\,L\Bigr),
\end{equation}
(con la convención de signos de Dorst y De Keninck~\cite{dorst-dekeninck}). El logaritmo recupera el eje, el ángulo y
el avance: es la manera natural de \emph{identificar el eje de un movimiento}.

\subsection{Velocidad: el tornillo instantáneo}
Si un cuerpo tiene motor $M(t)$, su velocidad es un bivector $\Omega$:
\begin{equation}
  \dot M = -\tfrac12\,M\,\Omega_{\text{c}} = -\tfrac12\,\Omega_{\text{e}}\,M,\qquad \Omega_{\text{e}} = M\,\Omega_{\text{c}}\,\widetilde{M},
\end{equation}
en el marco del cuerpo ($\Omega_{\text{c}}$) o del espacio ($\Omega_{\text{e}}$). $\Omega$ reúne velocidad angular y
lineal en un solo objeto, y su recta asociada es el \textbf{eje instantáneo de tornillo}: el simulador puede dibujar
alrededor de qué recta está girando cada pieza en cada instante.

\subsection{Alineación de agujeros, en PGA}
Para dos rectas normalizadas $L_1$, $L_2$, el producto $L_1L_2$ tiene parte escalar $-\cos\alpha$ y parte de grado 4
proporcional a $d\sin\alpha\,\mathbf{I}$, donde $\alpha$ es el ángulo entre ellas y $d$ su distancia mínima. Si son
paralelas, la distancia sale de la parte ideal de su conmutador. Así, la condición de alineación de la
sección~\ref{sec:ensamble} se reduce a medir $L_1L_2$.

\subsection{Dinámica en PGA}
La inercia es una aplicación lineal $\mathcal{A}$ de bivectores (velocidades) en bivectores duales (momentos): el
\emph{momento} $P = \mathcal{A}[\Omega]$ es una recta, la recta del momento, que reúne momento lineal y angular. Las
ecuaciones de Newton y Euler quedan en una sola, en el marco del cuerpo~\cite{dorst-dekeninck, gunn-tesis}:
\begin{equation}
  \dot\Omega_{\text{c}} = \mathcal{A}^{-1}\!\bigl[\,F_{\text{c}} - \Omega_{\text{c}}\times\mathcal{A}[\Omega_{\text{c}}]\,\bigr],
\end{equation}
donde $F$ es el \emph{forque} (fuerza y par juntos, también una recta: la recta de acción de la fuerza). La gravedad es
un forque vertical que pasa por el centro de masa. Sin forques externos, la recta del momento queda fija en el espacio:
es la conservación conjunta del momento lineal y angular, que se verifica en la prueba P13.

% ==========================================================================================
\section{Grafo de ensamble: tornillos, tuercas y uniones}\label{sec:ensamble}

\subsection{Definición}
El montaje es un grafo $G = (\mathcal{P}, \mathcal{U})$: los nodos son piezas (incluidos tornillos y tuercas) y las
aristas son uniones. Cada unión $u$ tiene un tipo $\tau(u)$ y un estado. El grafo es la «verdad» del montaje; los
cuerpos de la simulación se derivan de él.

\subsection{Alineación de agujeros}
Dos puertos $k$ (de la pieza $P$) y $l$ (de $Q$), llevados al mundo, están \emph{alineados} si sus ejes son paralelos y
coinciden:
\begin{equation}
  \lvert \vect{a}_k \times \vect{a}_l \rvert < \varepsilon_a, \qquad
  \bigl\lvert (\vect{o}_l - \vect{o}_k) - \bigl((\vect{o}_l-\vect{o}_k)\cdot\vect{a}_k\bigr)\vect{a}_k \bigr\rvert < \varepsilon_r,
\end{equation}
con $\varepsilon_a \approx \sin 3^\circ$ y $\varepsilon_r \approx 0{,}8$\,mm. Al soltar una pieza cerca de una
alineación, la interfaz la «imanta» a la posición exacta (ajuste a la retícula).

\subsection{El tornillo es una pieza}
Un tornillo $S$ se inserta a lo largo de una cadena de puertos alineados $k_1,\dots,k_n$ (en piezas
$P_1,\dots,P_n$), mientras su largo útil lo permita: $\sum_i \ell_{k_i} \le L_S$. Se forman las aristas
$(S,P_i)$ de tipo \emph{eje}. Según la tuerca:

\begin{tabularx}{\linewidth}{@{}l l Y@{}}
\toprule
\textbf{Estado} & \textbf{Unión de cada pieza con el tornillo} & \textbf{Grados de libertad relativos}\\
\midrule
Sin tuerca & cilíndrica & giro y deslizamiento a lo largo del eje (2): el tornillo puede salirse si nada lo
  detiene, y la gravedad puede sacarlo.\\
Tuerca floja & de revolución & giro alrededor del eje (1), con fricción pequeña.\\
Tuerca ajustada & de revolución con fricción & giro solo si el par supera $\tau_f$ (ec.~\ref{eq:tauf}).\\
Tuerca apretada & fija & ninguno: tornillo, tuerca y piezas forman un cuerpo.\\
Tuerca fija a una pieza & fija (tuerca–pieza) & la tuerca queda pegada a esa pieza (tuerca «cautiva»).\\
\bottomrule
\end{tabularx}

\medskip
Con tuerca floja, \emph{cada} pieza gira sobre el tornillo. Es lo que permite una biela articulada: dos vigas unidas
por un tornillo con tuerca floja.

\subsection{Apriete y fricción en la unión}
Apretar la tuerca con un par $T$ produce una fuerza de sujeción (precarga) $F_p$. Con la relación usual de
tornillería,
\begin{equation}
  T = K\, d\, F_p, \qquad K \approx 0{,}2,
\end{equation}
donde $d = 4$\,mm es el diámetro nominal. Las caras en contacto, una corona de radios $r_i$ y $r_o$ con presión
uniforme, resisten el giro hasta
\begin{equation}
  \tau_f = \mu_j\, F_p\, r_m, \qquad\label{eq:tauf} r_m = \frac{2}{3}\,\frac{r_o^3 - r_i^3}{r_o^2 - r_i^2}.
\end{equation}
En la simulación, la unión de revolución lleva un motor de velocidad cero con par máximo $\tau_f$: si el par que se le
aplica es menor, no gira; si es mayor, desliza con fricción cinética. Cuando $\tau_f$ supera cualquier par razonable
del montaje, la unión se convierte en \emph{fija} y las piezas se fusionan (sección~5). Esto se ve en pantalla: al
apretar la tuerca, el conjunto cambia de color de contorno y se mueve como un todo.

\subsection{Avance del tornillo}
Girar la tuerca un ángulo $\Delta\theta$ la desplaza a lo largo del eje
\begin{equation}
  \Delta s = p\,\frac{\Delta\theta}{2\pi}, \qquad p = 0{,}7\ \mathrm{mm}\ (\mathrm{M4}),
\end{equation}
que se usa para animar el apriete y para saber si la tuerca ya toca la pieza. El apriete solo crece cuando la tuerca
hace contacto.

% ==========================================================================================
\section{Cuerpos compuestos}

\subsection{Agrupación}
Las aristas \emph{fijas} definen una relación de equivalencia entre piezas. Sus clases (componentes conexas del
subgrafo de uniones fijas, calculadas con \emph{unión–búsqueda}) son los cuerpos rígidos de la simulación: «si las
junto, cambia el objeto como un todo». Cada vez que se aprieta, afloja, añade o quita algo, se recalculan solo las
clases afectadas.

\subsection{Propiedades de masa del compuesto}
Para un cuerpo $B$ formado por piezas $i$ con masas $m_i$, centros $\vect{c}_i$, inercias locales $\I_i$ y
orientaciones $\R_i$ (respecto al marco de $B$):
\begin{align}
  m_B &= \sum_i m_i, \qquad \vect{c}_B = \frac{1}{m_B}\sum_i m_i\,\vect{c}_i,\\
  \I_B &= \sum_i \Bigl[\R_i\,\I_i\,\R_i^{\mathsf T} + m_i\bigl(\lVert\vect{d}_i\rVert^2\E - \vect{d}_i\vect{d}_i^{\mathsf T}\bigr)\Bigr],
  \qquad \vect{d}_i = \vect{c}_i - \vect{c}_B,
\end{align}
que es el teorema de Steiner (ejes paralelos) en forma tensorial. Al fusionar, el cuerpo nuevo hereda el momento
lineal y angular total, para que fusionar no cree ni destruya movimiento:
\begin{equation}
  \vect{v}_B = \frac{1}{m_B}\sum_i m_i\vect{v}_i, \qquad
  \vect{\omega}_B = \I_B^{-1}\sum_i\Bigl[\I_i^{\text{mundo}}\vect{\omega}_i + m_i\,\vect{d}_i\times(\vect{v}_i - \vect{v}_B)\Bigr].
\end{equation}
La energía cinética sí puede bajar al fusionar (es un choque perfectamente inelástico); se registra como disipación.

\begin{decision}{Ventaja numérica}
Un tornillo de 2\,g unido a una rueda de 40\,g por una restricción es un sistema mal condicionado (cociente de masas
grande). Fusionarlos en un cuerpo elimina la restricción y el problema. Por eso fusionar no solo es fiel a la
realidad: también hace la simulación más estable.
\end{decision}

% ==========================================================================================
\section{Dinámica de cuerpo rígido}

Cada cuerpo obedece las ecuaciones de Newton–Euler en el marco del mundo:
\begin{align}
  m\,\dot{\vect{v}} &= m\vect{g} + \sum_k \vect{F}_k,\\
  \I\,\dot{\vect{\omega}} + \vect{\omega}\times\I\vect{\omega} &= \sum_k (\vect{p}_k - \vect{x})\times\vect{F}_k + \sum_j \vect{\tau}_j,
  \qquad \I = \R\,\I_{\text{local}}\R^{\mathsf T},
\end{align}
donde $\vect{F}_k$ son fuerzas aplicadas en puntos $\vect{p}_k$ (contactos, bandas, la mano) y $\vect{\tau}_j$
pares puros (motores, fricción en uniones). El término giroscópico $\vect{\omega}\times\I\vect{\omega}$ es el que hace
precesar un trompo hecho con una rueda.

\subsection{Restricciones (uniones)}
Una unión se escribe como una función $\vect{C}(\vect{q}) = \vect{0}$ sobre las configuraciones de los dos cuerpos.
Derivando, $\dot{\vect{C}} = \mathbf{J}\vect{u} = \vect{0}$, con $\vect{u} = (\vect{v}_a,\vect{\omega}_a,\vect{v}_b,\vect{\omega}_b)$
y $\mathbf{J}$ el jacobiano. Para la unión de revolución, con punto de anclaje $\vect{r}_a,\vect{r}_b$ y eje $\vect{e}$:
\begin{equation}
  \vect{C}_{\text{punto}} = \vect{x}_a + \R_a\vect{r}_a - \vect{x}_b - \R_b\vect{r}_b = \vect{0}\ (3\ \text{ec.}),\qquad
  (\R_a\vect{e})\cdot(\R_b\vect{n}_1) = (\R_a\vect{e})\cdot(\R_b\vect{n}_2) = 0\ (2\ \text{ec.}),
\end{equation}
donde $\vect{n}_1,\vect{n}_2$ son perpendiculares al eje en $b$. La unión cilíndrica quita una ecuación del punto
(la del eje) y la fija las impone todas (6). El núcleo resuelve los impulsos $\vect{\lambda}$ con
\begin{equation}
  \mathbf{J}\mathbf{M}^{-1}\mathbf{J}^{\mathsf T}\,\vect{\lambda} = -\mathbf{J}\vect{u} - \frac{\beta}{\Delta t}\vect{C},
\end{equation}
por Gauss–Seidel proyectado, con $\beta$ la corrección de deriva (estabilización de Baumgarte).

\subsection{Holgura}
El agujero (4{,}6\,mm) es más ancho que el tornillo (4{,}0\,mm): hay $0{,}3$\,mm de juego radial. En la fase 1 se
ignora (unión ideal). Si se quiere, se modela después con un límite de distancia en lugar de una unión exacta, para
que las piezas «bailen» como en la realidad.

% ==========================================================================================
\section{Contacto y fricción}

\subsection{Ley de Coulomb}
En cada punto de contacto con normal $\vect{n}$, la fuerza se descompone en normal $F_n\ge 0$ y tangencial
$\vect{F}_t$, restringida al cono de fricción:
\begin{equation}
  \lVert\vect{F}_t\rVert \le \mu_s F_n \quad\text{(adherencia)}, \qquad
  \vect{F}_t = -\mu_k F_n\,\frac{\vect{v}_t}{\lVert\vect{v}_t\rVert} \quad\text{(deslizamiento)},
\end{equation}
con $\vect{v}_t$ la velocidad relativa tangencial. La restitución $e$ fija el rebote:
$v_n^{+} = -e\,v_n^{-}$.

\subsection{Materiales}
Entre dos materiales se combinan los coeficientes con la media geométrica, $\mu_{12} = \sqrt{\mu_1\mu_2}$, salvo que la
tabla dé el par directamente.

\begin{tabularx}{\linewidth}{@{}Y c c c@{}}
\toprule
\textbf{Par de materiales} & $\mu_s$ & $\mu_k$ & $e$\\
\midrule
plástico–plástico & 0,35 & 0,30 & 0,4\\
metal–plástico & 0,30 & 0,25 & 0,3\\
caucho (aro, banda)–plástico & 0,80 & 0,70 & 0,6\\
cualquiera–mesa (madera) & 0,45 & 0,40 & 0,3\\
\bottomrule
\end{tabularx}

\smallskip
{\small Valores iniciales típicos; se reemplazan por los medidos (sección~\ref{sec:calibracion}).}

\subsection{Detección y solución del contacto}
Cada pieza tiene una forma de colisión convexa o unión de convexas (cajas, cilindros, cápsulas). La detección se hace
en dos etapas: una gruesa con cajas alineadas a los ejes, y una fina con el teorema del eje separador (SAT) o con
GJK y EPA~\cite{ericson, vandenbergen}, que dan la normal, la profundidad y los puntos de contacto. El contacto es una
ligadura \emph{unilateral} (condición de Signorini: $F_n\ge0$, separación $\ge0$, y su producto nulo), y con la
fricción de Coulomb forma un problema de complementariedad que se resuelve iterativamente (Gauss–Seidel proyectado,
con impulsos acumulados~\cite{catto, stewart-trinkle}).

\subsection{Fricción de rodadura}
Una rueda que rueda sobre la mesa pierde energía por deformación. Se modela con un par opuesto al giro
$\tau_r = c_r F_n R$, con $c_r \approx 0{,}01$, para que una rueda lanzada termine deteniéndose.

% ==========================================================================================
\section{Ruedas de fricción}

Dos ruedas que se tocan transmiten movimiento por fricción. La colisión genérica entre cilindros sirve
para casos sueltos, pero para que la relación de transmisión sea exacta y explicable se usa un modelo propio que se
activa cuando la capa de ensamble detecta contacto entre dos discos.

\subsection{Condición general de rodadura}
Sean dos cuerpos $A$ y $B$ que giran alrededor de ejes fijos (a su estructura) con direcciones $\vect{a}$, $\vect{b}$,
centros $\vect{c}_A$, $\vect{c}_B$ y velocidades angulares $\omega_A\vect{a}$, $\omega_B\vect{b}$. En el punto de
contacto $\vect{P}$, con tangente común $\vect{t}$, las velocidades de cada superficie son
\begin{equation}
  \vect{v}_A(\vect{P}) = \omega_A\,\vect{a}\times(\vect{P}-\vect{c}_A), \qquad
  \vect{v}_B(\vect{P}) = \omega_B\,\vect{b}\times(\vect{P}-\vect{c}_B).
\end{equation}
Rodar sin deslizar significa que la velocidad relativa tangencial es nula:
\begin{equation}
  s \equiv \bigl(\vect{v}_A(\vect{P}) - \vect{v}_B(\vect{P})\bigr)\cdot\vect{t} = 0 .
\end{equation}

\subsection{Caso paralelo (canto con canto)}
Ejes paralelos, ruedas de radios $R_A$, $R_B$ tocándose por el borde:
\begin{equation}
  \omega_A R_A = -\,\omega_B R_B \quad\Longrightarrow\quad i = \frac{\omega_B}{\omega_A} = -\frac{R_A}{R_B}.
\end{equation}
El signo menos indica que giran en sentidos opuestos (el sistema de dos ruedas de la cartilla, con $i=-1/3$).

\subsection{Caso transversal (canto contra cara)}
Una rueda $A$ de radio $R$ apoya su canto sobre la \emph{cara} de un disco $B$, con ejes perpendiculares, a una
distancia $r$ del eje de $B$. Ahora la velocidad de la superficie de $B$ en el contacto es $\omega_B r$, y
\begin{equation}
  \boxed{\;\omega_A R = \omega_B\, r \quad\Longrightarrow\quad \frac{\omega_A}{\omega_B} = \frac{r}{R}\;}
\end{equation}
La relación \textbf{depende de dónde está el contacto}: si se corre la rueda pequeña hacia el borde del disco, gira
más rápido; si se acerca al centro, más lento; y al pasar al otro lado del centro, invierte el sentido. Es una
transmisión de relación variable continua (como el variador de disco y rueda). El simulador mostrará $r$ y la
relación en vivo mientras se desliza la rueda.

\subsection{Deslizamiento}
El contacto transmite como máximo la fuerza tangencial $\mu_s N$, donde $N$ es la fuerza normal (del peso, de una
banda que empuja la rueda contra el disco o de la sujeción). Entonces el par máximo sobre el disco es
\begin{equation}
  \tau_{B,\max} = \mu_s N\, r .
\end{equation}
Si se pide más, la rueda patina, y la fuerza pasa a ser cinética. Para evitar la discontinuidad del signo se usa
una versión regularizada:
\begin{equation}
  F_t = -\mu_k N \tanh\!\left(\frac{s}{s_0}\right),\qquad s_0 \approx 0{,}5\ \mathrm{mm/s},
\end{equation}
aplicada como fuerzas iguales y opuestas en $\vect{P}$ sobre $A$ y $B$ (tercera ley), que producen pares
$\vect{\tau}_A = (\vect{P}-\vect{c}_A)\times F_t\vect{t}$ y $\vect{\tau}_B = -(\vect{P}-\vect{c}_B)\times F_t\vect{t}$.

\subsection{Pérdida por pivoteo}
En el caso transversal, el contacto no es un punto sino una franja de ancho $w$ (el grosor de la rueda) a lo largo
del radio del disco. Sus extremos, en $r\pm w/2$, se mueven a $\omega_B(r\pm w/2)$ pero la rueda solo puede igualar
una velocidad. Con presión uniforme, el deslizamiento medio es $\overline{\lvert\Delta v\rvert} = \omega_B w/4$, y se
disipa una potencia
\begin{equation}
  P_{\text{pivoteo}} \approx \mu_k N\,\frac{\omega_B\, w}{4},
\end{equation}
que explica por qué esta transmisión calienta y desgasta más que la de canto con canto, y por qué conviene una rueda
delgada.

% ==========================================================================================
\section{Bandas elásticas}

\subsection{Elasticidad del caucho}
Una banda de longitud natural $L_0$, estirada hasta $L$, ejerce una tensión
\begin{equation}
  T = \begin{cases} k\,(L - L_0) + c\,\dot L, & L > L_0,\\ 0, & L \le L_0 \ \text{(floja)},\end{cases}
  \qquad k = \frac{E A}{L_0},
\end{equation}
con $E$ el módulo de Young del caucho ($\sim 1$--$3$\,MPa), $A$ el área de su sección y $c$ un amortiguamiento
pequeño (el caucho disipa al estirarse y recogerse: histéresis). Para estiramientos grandes ($L > 2L_0$) el caucho
se endurece; si hace falta, se usa la ley neo-Hookeana $T = \tfrac{EA}{3}\bigl(\lambda - \lambda^{-2}\bigr)$, con
$\lambda = L/L_0$. La constante $k$ de cada banda se mide colgándole pesas (sección~\ref{sec:calibracion}).

\subsection{Banda libre entre dos puntos}
Si la banda une dos puntos $\vect{p}_1$, $\vect{p}_2$ (un gancho y un tornillo, una honda), aplica fuerzas iguales y
opuestas a lo largo de la recta:
\begin{equation}
  \vect{F}_1 = T\,\frac{\vect{p}_2 - \vect{p}_1}{\lVert\vect{p}_2 - \vect{p}_1\rVert} = -\vect{F}_2,
  \qquad L = \lVert\vect{p}_2 - \vect{p}_1\rVert .
\end{equation}

\subsection{Banda sobre dos poleas}
Poleas de radios $R_1 \ge R_2$ con centros a distancia $D$. Con la banda abierta (sin cruzar), la geometría da
\begin{equation}
  \beta = \arcsin\frac{R_1 - R_2}{D}, \qquad
  L = 2\sqrt{D^2 - (R_1-R_2)^2} + R_1(\pi + 2\beta) + R_2(\pi - 2\beta),
\end{equation}
y ángulos abrazados $\theta_1 = \pi + 2\beta$ y $\theta_2 = \pi - 2\beta$. Con la banda cruzada (las poleas giran en
sentidos opuestos):
\begin{equation}
  \beta_\times = \arcsin\frac{R_1 + R_2}{D}, \qquad
  L_\times = 2\sqrt{D^2 - (R_1+R_2)^2} + (R_1 + R_2)(\pi + 2\beta_\times).
\end{equation}
Para más de dos poleas, $L$ es la suma de tramos tangentes y arcos abrazados, calculados en el plano de las poleas.
La precarga es $T_0 = k(L - L_0)$: mover las poleas para alejarlas tensa la banda.

\subsection{Transmisión y deslizamiento}
Sin deslizamiento, $\omega_1 R_1 = \omega_2 R_2$ (abierta) u $\omega_1 R_1 = -\omega_2 R_2$ (cruzada). Al transmitir un
par $\tau$, un tramo se tensa y el otro se afloja. Si la longitud total se conserva y el caucho es lineal,
\begin{equation}
  T_1 = T_0 + \frac{\tau}{2R}, \qquad T_2 = T_0 - \frac{\tau}{2R}.
\end{equation}
La banda no resbala mientras se cumpla la ecuación de Euler–Eytelwein (del cabrestante) en la polea con menor ángulo
abrazado:
\begin{equation}
  \frac{T_1}{T_2} \le e^{\mu\theta},
\end{equation}
de donde sale el par máximo que la banda puede transmitir:
\begin{equation}
  \boxed{\;\tau_{\max} = 2\,T_0\,R\;\frac{e^{\mu\theta}-1}{e^{\mu\theta}+1}\;}\label{eq:taumax}
\end{equation}
Más tensión inicial o más ángulo abrazado permiten más par; por eso una banda floja patina.

\subsection{Modelo dinámico}
La banda se representa como un acople elástico entre los ángulos $\varphi_1$, $\varphi_2$ de las poleas, con la
diferencia de tensiones
\begin{equation}
  T_1 - T_2 = k_e\,\bigl(R_1\varphi_1 - R_2\varphi_2 - \delta\bigr) + c_e\,\bigl(R_1\omega_1 - R_2\omega_2\bigr),
  \qquad k_e = \frac{2k L_0}{L_{\text{tramo}}},
\end{equation}
donde $\delta$ es el deslizamiento acumulado. Si $\lvert T_1 - T_2\rvert$ supera el límite del cabrestante, se satura en
ese valor y $\delta$ avanza: la banda patina. Cada polea recibe el par $\mp R_i(T_1 - T_2)$ y su eje la fuerza
$(T_1 + T_2)$ hacia la otra polea. Así se ven, con el mismo modelo, la transmisión normal, el arranque elástico y el
deslizamiento.

% ==========================================================================================
\section{Análisis del mecanismo: mecánica analítica}\label{sec:analitica}

La simulación avanza con ligaduras resueltas por impulsos. La mecánica analítica se usa en paralelo para
\emph{entender, contar y explicar} el mecanismo armado, y como verificación.

\subsection{Configuración y ligaduras}
Un montaje con $N$ cuerpos tiene configuración $(M_1,\dots,M_N)$, un punto de $SE(3)^N$, de dimensión $6N$. Las
uniones imponen ligaduras $\vect{C}(\vect{q}) = \vect{0}$, y sus velocidades admisibles cumplen
$\mathbf{J}\,\vect{u} = \vect{0}$.
\begin{itemize}[leftmargin=*]
  \item \textbf{Holónomas}: se pueden escribir como relación entre posiciones (uniones, dos ruedas con ejes fijos:
        $R_A\varphi_A + R_B\varphi_B = \text{cte}$, que es integrable).
  \item \textbf{No holónomas}: solo entre velocidades, no integrables (una rueda que rueda libre sobre la mesa).
\end{itemize}

\subsection{Grados de libertad y ejes}
El número de grados de libertad es
\begin{equation}
  \text{GDL} = 6N - \operatorname{rango}\mathbf{J},
\end{equation}
calculado numéricamente con una descomposición en valores singulares de $\mathbf{J}$. La fórmula de
Grübler–Kutzbach, $\text{GDL} = 6(n - 1 - j) + \sum_i f_i$ (con $n$ eslabones incluido el fijo, $j$ uniones y $f_i$
los grados de libertad de cada una), sirve como contraste: cuando no coinciden, hay \emph{ligaduras redundantes} (el
montaje está sobrerrestringido), y el simulador lo dice. Los vectores del núcleo de $\mathbf{J}$ son los movimientos
posibles; cada uno es un bivector por cuerpo, es decir, un \textbf{eje de giro} que se puede dibujar. En pantalla:
«este montaje tiene 1 grado de libertad: todo depende del ángulo de la manivela», o «0: está trabado».

\subsection{Coordenadas generalizadas y lagrangiano reducido}
En un mecanismo de un grado de libertad, todo depende de una coordenada $q$ (el ángulo de la manivela). Con
$\Omega_i = \vect{\xi}_i(q)\,\dot q$, la energía cinética y la potencial son
\begin{equation}
  T = \tfrac12\,\mathcal{M}(q)\,\dot q^{\,2},\qquad \mathcal{M}(q) = \sum_i \xi_i(q)\vee\mathcal{A}_i[\xi_i(q)],\qquad
  V(q) = \sum_i m_i\,g\,y_i(q),
\end{equation}
y la ecuación de Euler–Lagrange da la ecuación de movimiento del mecanismo entero:
\begin{equation}
  \mathcal{M}(q)\,\ddot q + \tfrac12\,\mathcal{M}'(q)\,\dot q^{\,2} + V'(q) = Q,
\end{equation}
con $Q$ la fuerza generalizada (el par en la manivela). $\mathcal{M}(q)$ es la \emph{inercia efectiva}: explica por
qué la biela-manivela se siente «más pesada» en ciertos ángulos.

\subsection{Trabajo virtual y ventaja mecánica}
En equilibrio, el trabajo virtual de las fuerzas aplicadas es nulo para todo desplazamiento admisible:
$\sum_k F_k\vee\delta\Lambda_k = 0$. De ahí salen, sin dibujar fuerzas internas, el equilibrio de la balanza
($m_1x_1 = m_2x_2$) y la ventaja mecánica de cualquier transmisión,
$\tau_{\text{sal}}/\tau_{\text{ent}} = \delta\theta_{\text{ent}}/\delta\theta_{\text{sal}}$.

\subsection{Fuerzas de ligadura}
Los multiplicadores de Lagrange $\vect{\lambda}$ que calcula el núcleo son las fuerzas que soporta cada unión. En PGA
cada una es un forque (una recta de fuerza). El simulador puede mostrarlas como flechas: «este tornillo soporta
0{,}3\,N».

% ==========================================================================================
\section{Interacción}\label{sec:interaccion}

\begin{itemize}[leftmargin=*]
  \item \textbf{Arrastrar una pieza}: la mano es un resorte amortiguado entre el cursor $\vect{m}$ y el punto tomado
        $\vect{p}$: $\vect{F} = k_m(\vect{m} - \vect{p}) - c_m\dot{\vect{p}}$, con $c_m = 2\sqrt{k_m m}$ (amortiguamiento
        crítico) y fuerza limitada, para que la mano no atraviese ni rompa nada.
  \item \textbf{Manivela}: un motor en la unión de revolución que sigue la velocidad que pide el dedo, con par máximo
        $\tau_{\text{mano}}$. Si el mecanismo se traba, la manivela no gira; no se fuerza.
  \item \textbf{Herramientas}: llave para apretar o aflojar tuercas (sección~4), banda elástica (se tiende tocando
        poleas o puntos), regla para medir, cronómetro.
  \item \textbf{Modos}: \emph{armar} (física en pausa, piezas imantadas a la retícula) y \emph{simular} (física
        activa). Se puede pausar, avanzar cuadro a cuadro y ver en cámara lenta.
\end{itemize}

% ==========================================================================================
\section{Integración numérica y estabilidad}\label{sec:integracion}

\begin{itemize}[leftmargin=*]
  \item \textbf{Paso}: el dibujo va a 60 cuadros por segundo; la física avanza con $\Delta t = 1/60$\,s dividido en
        $n = 4$ a $8$ subpasos, es decir, $\Delta t_{\text{sub}} \approx 2$--$4$\,ms. Se avanza con Euler semiimplícito
        sobre el grupo: primero la velocidad, $\Omega_{k+1} = \Omega_k + \Delta t\,\dot\Omega_k$, y luego el motor con
        la exponencial, $M_{k+1} = M_k\exp\!\bigl(-\tfrac{\Delta t}{2}\Omega_{k+1}\bigr)$, que mantiene $M$ como
        movimiento rígido exacto (sin renormalizar)~\cite{hairer, iserles}.
  \item \textbf{Resortes} (bandas, mano): un resorte explícito de rigidez $k$ sobre una masa efectiva $m$ es estable si
        $\omega_n\,\Delta t_{\text{sub}} < 2$, con $\omega_n = \sqrt{k/m}$. Las bandas muy rígidas o muy livianas se
        tratan de forma implícita, o con la formulación XPBD de rigidez compatible~\cite{xpbd}.
  \item \textbf{Velocidades altas}: detección continua de colisiones para piezas pequeñas y rápidas (tornillos que caen).
  \item \textbf{Reposo}: los cuerpos quietos «duermen» para ahorrar cálculo, y despiertan al tocarlos.
  \item \textbf{Energía}: se lleva la cuenta $E = \sum(\tfrac12 m v^2 + \tfrac12\vect{\omega}\cdot\I\vect{\omega} + m g y)
        + \sum\tfrac12 k (L-L_0)_+^2$ y lo disipado por fricción. Sin fricción, $E$ debe conservarse (prueba P2).
\end{itemize}

% ==========================================================================================
\section{Calibración con el kit real}\label{sec:calibracion}

\begin{tabularx}{\linewidth}{@{}l Y l@{}}
\toprule
\textbf{Parámetro} & \textbf{Experimento} & \textbf{Fórmula}\\
\midrule
Masa de cada pieza & Pesar 10 piezas iguales en una báscula de cocina y dividir. & $m = M/10$\\
Fricción estática & Plano inclinado: subir el ángulo hasta que la pieza empieza a deslizar. & $\mu_s = \tan\theta_c$\\
Fricción cinética & Plano más inclinado: medir el tiempo $t$ en recorrer $d$ desde el reposo. &
  $\mu_k = \tan\theta - \dfrac{2d}{g t^2\cos\theta}$\\
Rigidez de banda & Colgar pesas $m_j$ y medir el estiramiento $x_j$; ajustar una recta. & $k = \dfrac{\Delta(m g)}{\Delta x}$\\
Fricción en el eje & Péndulo de viga sobre tornillo: contar oscilaciones hasta detenerse. & decaimiento lineal de la amplitud\\
\bottomrule
\end{tabularx}

\smallskip
Estos experimentos son, además, actividades para el aula: los estudiantes miden el kit y alimentan el simulador.

% ==========================================================================================
\section{Verificación: pruebas con respuesta conocida}\label{sec:verificacion}

Cada fase se acepta solo si sus pruebas pasan, con un error relativo menor al indicado.

\begin{tabularx}{\linewidth}{@{}l Y l c@{}}
\toprule
& \textbf{Montaje} & \textbf{Resultado analítico} & \textbf{Tolerancia}\\
\midrule
P1 & Tornillo en caída libre durante 0,5\,s & $y = -\tfrac12 g t^2$ & 1\,\%\\
P2 & Viga colgada de un tornillo, tuerca floja, sin fricción (péndulo físico) &
  $T = 2\pi\sqrt{I_O/(m g d)}$; $E$ constante & 2\,\%\\
P3 & Pieza sobre plano inclinado & desliza solo si $\tan\theta > \mu_s$ & exacto\\
P4 & Balanza de brazos iguales con pesas & equilibrio si $m_1 x_1 = m_2 x_2$ & 2\,\%\\
P5 & Dos ruedas de canto, $R_A/R_B = 1/3$ & $\omega_B/\omega_A = -1/3$ & 1\,\%\\
P6 & Rueda sobre la cara de un disco, a radio $r$ & $\omega_A/\omega_B = r/R$ & 2\,\%\\
P7 & Banda sobre dos poleas, par creciente & patina al llegar a $\tau_{\max}$ (ec.~\ref{eq:taumax}) & 5\,\%\\
P8 & Apretar la tuerca de dos vigas en movimiento & se conserva el momento angular total & 1\,\%\\
P9 & Tornillo sin tuerca, eje vertical, pieza invertida & el tornillo se cae por gravedad & cualitativa\\
P10 & Identidades de PGA: $M\widetilde M = 1$, $\exp(\log M) = M$, composición, sándwich de puntos y rectas & exactas & $10^{-12}$\\
P11 & Cuerpo libre sin gravedad, con giro inicial, 10\,s & recta del momento fija; energía constante & 0,1\,\%\\
P12 & Grados de libertad: balanza, biela-manivela, montaje trabado & 1, 1 y 0 & exacto\\
P13 & Balanza por trabajo virtual y por simulación & mismo equilibrio & 1\,\%\\
P14 & Lagrangiano reducido frente a simulación completa (biela-manivela) & misma $q(t)$ & 2\,\%\\
\bottomrule
\end{tabularx}

\smallskip
Las pruebas corren solas (sin pantalla) en el navegador sin interfaz, igual que las pruebas actuales de la web, y se
repiten tras cada cambio.

% ==========================================================================================
\section{Módulos de estudio y construcción}\label{sec:modulos}

Cada módulo tiene su ficha de estudio, sus pruebas y su función pública.

\medskip
\noindent\begin{tabularx}{\linewidth}{@{}p{2.1cm} Y Y p{2.3cm}@{}}
\toprule
\textbf{Módulo} & \textbf{Estudio} & \textbf{Construcción} & \textbf{Pruebas}\\
\midrule
0. Álgebra & Álgebra geométrica: producto geométrico, grados, rotores; luego PGA: planos, rectas, puntos, motores,
  exponencial y logaritmo~\cite{hestenes-nfcm, doran-lasenby, gunn-dekeninck, dorst-dekeninck, lengyel}. &
  \texttt{pga.js}; adaptador motor $\to$ matriz de three.js; rayo del cursor como recta; mesa de trabajo. & P10\\
1. Cinemática & Movimientos rígidos, Chasles, tornillos, velocidad como bivector, teoría de tornillos
  clásica~\cite{ball, selig, murray, lynch-park}. & Piezas con puertos (rectas), mover, girar, imantar a la retícula,
  alineación de agujeros. & P10, alineación\\
2. Dinámica & Newton–Euler, tensor de inercia, Steiner, cuerpo rígido en PGA, integración en grupos de
  Lie~\cite{goldstein, dorst-dekeninck, gunn-tesis, hairer, iserles}. & \texttt{cuerpo.js}, \texttt{integrador.js};
  masa e inercia de las piezas; gravedad; caída y giro libre (sin choques). & P1, P11\\
3. Mecánica analítica I & Ligaduras, coordenadas generalizadas, multiplicadores, holonomía, Grübler, rango del
  jacobiano~\cite{lanczos, goldstein, arnold, lynch-park}. & Jacobianos de las uniones en PGA; contador de grados de
  libertad; dibujar los ejes posibles. & P12\\
4. Ligaduras & Ligaduras dinámicas, potencia virtual en PGA, estabilización, Gauss–Seidel proyectado,
  XPBD~\cite{hadfield-lasenby, baraff, catto, xpbd, muller-rigid}. & Tornillos, tuercas y sus estados; cuerpos
  compuestos; fusión con conservación del momento. & P2, P8, P9\\
5. Contacto & Detección (SAT, GJK, EPA), condición de Signorini, cono de Coulomb, complementariedad, apilamiento
  estable~\cite{ericson, vandenbergen, stewart-trinkle, catto}. & \texttt{colision.js}, \texttt{contacto.js}; la mesa,
  choques, apilar piezas. & P3\\
6. Transmisiones & Rodadura, ruedas de fricción, mecánica de contacto con deslizamiento y pivoteo~\cite{johnson},
  bandas y cabrestante, elasticidad del caucho~\cite{treloar}. & Ruedas paralelas y transversales; bandas con
  deslizamiento; manivela. & P5, P6, P7\\
7. Mecánica analítica II & Lagrangiano reducido, inercia efectiva, trabajo virtual, fuerzas de ligadura,
  integradores variacionales~\cite{lanczos, marsden-west, doran-lasenby}. & «Ver las cuentas»: la ecuación del
  mecanismo armado, fuerzas en cada tornillo, ventaja mecánica. & P4, P13, P14\\
8. Compartir & --- & Guardar y cargar montajes, retos con los 7 prototipos, exportar a realidad aumentada. & ---\\
\bottomrule
\end{tabularx}

\begin{decision}{Ritmo}
Un módulo no se cierra hasta que su ficha está estudiada, sus pruebas pasan y su función pública funciona. Los
módulos 0 a 2 dan la base (álgebra, cinemática, dinámica); el 5 (contacto) es el más largo y el más difícil.
\end{decision}

% ==========================================================================================
\section{Riesgos y decisiones abiertas}

\begin{tabularx}{\linewidth}{@{}Y Y@{}}
\toprule
\textbf{Riesgo} & \textbf{Cómo se maneja}\\
\midrule
Construir el motor desde cero toma mucho tiempo & Módulos con entregas visibles desde el módulo 0; respaldo con una
  librería de colisiones solo si el módulo 5 se atasca.\\
Convenciones de signo en PGA (hay varias en la literatura) & Se fija una (la de Dorst y De Keninck) en la ficha 0 y
  se verifica con pruebas antes de usarla.\\
Rendimiento en celulares con muchos cuerpos & Fusión en cuerpos compuestos, formas de colisión simples, reposo
  automático; límite de piezas en celular.\\
Inestabilidad con piezas muy livianas o bandas rígidas & Escala CGS, subpasos, resortes implícitos (sección~\ref{sec:integracion}).\\
Armar en 3D con el dedo es difícil & Imantado a la retícula, giros de 90° con un toque, vista ortogonal opcional.\\
Parámetros físicos inventados & Calibración con el kit (sección~\ref{sec:calibracion}); todo parámetro visible y editable.\\
\bottomrule
\end{tabularx}

\medskip
\textbf{Decisiones abiertas:}
\begin{enumerate}[leftmargin=*]
  \item Nombre del programa (propuesta: \emph{Taller virtual ATA EPE}).
  \item Si la cara pública vive dentro de la web de la cartilla (\texttt{/taller/}) o en el futuro sitio de ATA EPE
        Robótica Escolar. Por ahora se desarrolla en \texttt{/taller/}.
  \item Qué piezas faltan en el catálogo para los prototipos: aros de caucho, poleas, manivela, ganchos, bandas.
\end{enumerate}

% ==========================================================================================
\begin{thebibliography}{99}
\small

\item[] \textbf{Álgebra geométrica y su aplicación a la mecánica}
\bibitem{hestenes-nfcm} D.~Hestenes, \emph{New Foundations for Classical Mechanics}, 2.ª ed., Kluwer, 1999.
  La reformulación completa de la mecánica clásica en álgebra geométrica.
\bibitem{doran-lasenby} C.~Doran y A.~Lasenby, \emph{Geometric Algebra for Physicists}, Cambridge University
  Press, 2003. Caps.~2--4 (álgebra, mecánica clásica, cuerpo rígido) y 12 (lagrangianos con multivectores).
\bibitem{dorst-fontijne-mann} L.~Dorst, D.~Fontijne y S.~Mann, \emph{Geometric Algebra for Computer Science},
  Morgan Kaufmann, 2007.
\bibitem{lengyel} E.~Lengyel, \emph{Projective Geometric Algebra Illuminated}, Terathon Software, 2024.
\bibitem{gunn-tesis} C.~Gunn, \emph{Geometry, Kinematics, and Rigid Body Mechanics in Cayley-Klein Geometries},
  tesis doctoral, TU Berlín, 2011. doi:10.14279/depositonce-3058. Cap.~9: dinámica del cuerpo rígido en PGA.
\bibitem{gunn-dekeninck} C.~Gunn y S.~De Keninck, \emph{Geometric Algebra for Computer Graphics}, notas del curso,
  SIGGRAPH 2019.
\bibitem{dorst-dekeninck} L.~Dorst y S.~De Keninck, \emph{May the Forque Be with You: Dynamics in PGA},
  v2.6, 2023. \url{https://bivector.net/PGADYN.html}
\bibitem{hadfield-lasenby} H.~Hadfield y J.~Lasenby, «Constrained dynamics in conformal and projective geometric
  algebra», en \emph{Advances in Computer Graphics (CGI 2020)}, LNCS, Springer, 2020, pp.~459--471.
\bibitem{hestenes-fasse} D.~Hestenes y E.~D.~Fasse, «Homogeneous rigid body mechanics with elastic coupling», en
  L.~Dorst, C.~Doran y J.~Lasenby (eds.), \emph{Applications of Geometric Algebra in Computer Science and
  Engineering}, Birkhäuser, 2002. doi:10.1007/978-1-4612-0089-5\_19
\bibitem{clifford} W.~K.~Clifford, «Preliminary sketch of biquaternions», \emph{Proc.\ London Math.\ Soc.} 4 (1873)
  381--395. El origen histórico de los motores.

\item[] \textbf{Teoría de tornillos y cinemática de mecanismos}
\bibitem{ball} R.~S.~Ball, \emph{A Treatise on the Theory of Screws}, Cambridge University Press, 1900.
\bibitem{selig} J.~M.~Selig, \emph{Geometric Fundamentals of Robotics}, 2.ª ed., Springer, 2005.
\bibitem{murray} R.~M.~Murray, Z.~Li y S.~S.~Sastry, \emph{A Mathematical Introduction to Robotic Manipulation},
  CRC Press, 1994.
\bibitem{lynch-park} K.~M.~Lynch y F.~C.~Park, \emph{Modern Robotics: Mechanics, Planning, and Control},
  Cambridge University Press, 2017. Cap.~2 (grados de libertad, Grübler) y cap.~3 (tornillos).
\bibitem{mirtich} B.~Mirtich, «Fast and accurate computation of polyhedral mass properties», \emph{Journal of
  Graphics Tools} 1(2) (1996) 31--50.
\bibitem{featherstone} R.~Featherstone, \emph{Rigid Body Dynamics Algorithms}, Springer, 2008.

\item[] \textbf{Mecánica analítica}
\bibitem{goldstein} H.~Goldstein, C.~Poole y J.~Safko, \emph{Classical Mechanics}, 3.ª ed., Addison-Wesley, 2002.
\bibitem{lanczos} C.~Lanczos, \emph{The Variational Principles of Mechanics}, 4.ª ed., Dover, 1986.
\bibitem{arnold} V.~I.~Arnold, \emph{Mathematical Methods of Classical Mechanics}, 2.ª ed., Springer, 1989.
\bibitem{landau} L.~D.~Landau y E.~M.~Lifshitz, \emph{Mecánica}, Reverté.

\item[] \textbf{Integración numérica geométrica}
\bibitem{hairer} E.~Hairer, C.~Lubich y G.~Wanner, \emph{Geometric Numerical Integration}, 2.ª ed., Springer, 2006.
\bibitem{iserles} A.~Iserles, H.~Munthe-Kaas, S.~Nørsett y A.~Zanna, «Lie-group methods», \emph{Acta Numerica} 9
  (2000) 215--365.
\bibitem{marsden-west} J.~E.~Marsden y M.~West, «Discrete mechanics and variational integrators», \emph{Acta
  Numerica} 10 (2001) 357--514.

\item[] \textbf{Ligaduras, contacto y colisiones}
\bibitem{baraff} D.~Baraff y A.~Witkin, \emph{Physically Based Modeling}, notas del curso, SIGGRAPH 2001.
\bibitem{catto} E.~Catto, «Iterative dynamics with temporal coherence», Game Developers Conference, 2005.
\bibitem{stewart-trinkle} D.~E.~Stewart y J.~C.~Trinkle, «An implicit time-stepping scheme for rigid body dynamics
  with inelastic collisions and Coulomb friction», \emph{Int.\ J.\ Numer.\ Methods Eng.} 39 (1996) 2673--2691.
\bibitem{xpbd} M.~Macklin, M.~Müller y N.~Chentanez, «XPBD: position-based simulation of compliant constrained
  dynamics», \emph{Motion in Games}, 2016.
\bibitem{muller-rigid} M.~Müller et al., «Detailed rigid body simulation with extended position based dynamics»,
  \emph{Computer Graphics Forum} 39(8), 2020.
\bibitem{ericson} C.~Ericson, \emph{Real-Time Collision Detection}, Morgan Kaufmann, 2004.
\bibitem{vandenbergen} G.~van den Bergen, \emph{Collision Detection in Interactive 3D Environments}, Morgan
  Kaufmann, 2003.

\item[] \textbf{Contacto, fricción, materiales y transmisiones}
\bibitem{shigley} R.~G.~Budynas y J.~K.~Nisbett, \emph{Shigley's Mechanical Engineering Design}, 10.ª ed.,
  McGraw-Hill, 2015. Cap.~17: transmisiones por banda.
\bibitem{meriam} J.~L.~Meriam y L.~G.~Kraige, \emph{Engineering Mechanics: Statics}, Wiley. Fricción en bandas.
\bibitem{johnson} K.~L.~Johnson, \emph{Contact Mechanics}, Cambridge University Press, 1985. Rodadura con
  deslizamiento y pivoteo.
\bibitem{treloar} L.~R.~G.~Treloar, \emph{The Physics of Rubber Elasticity}, 3.ª ed., Oxford, 1975.

\item[] \textbf{Formatos y computación}
\bibitem{gltf} Khronos Group, \emph{glTF 2.0 Specification}. \url{https://registry.khronos.org/glTF/}
\bibitem{goldberg} D.~Goldberg, «What every computer scientist should know about floating-point arithmetic», \emph{ACM
  Computing Surveys} 23(1) (1991) 5--48.
\end{thebibliography}

\end{document}
